\documentclass{article}
\usepackage{graphicx} 
\usepackage{amsmath}
\usepackage{amssymb}

\title{Trabajo MAC: NP-Completitud}
\author{RAMÓN DOMECH LEÓN}
\date{Junio 2026}

\begin{document}
\setlength{\parindent}{0pt}
\maketitle

\section*{Demostración de la NP-completitud del problema 10.h Partición en Tríos}

\subsection*{Definición del problema}

Llamaremos \textbf{Partición en Tríos (PT)} al siguiente problema de decisión: Dado un conjunto de enteros cuya longitud es un múltiplo de
3, ¿es posible dividir el conjunto en tríos (conjuntos de 3 elementos) de forma que
los enteros de cada trío sumen lo mismo?

\vspace{1\baselineskip}
Reescribamos el enunciado diferenciando qué constituye una instancia del problema y cuál es la pregunta exacta a resolver:

\medskip

\textbf{Instancia:} Un conjunto $A$ de enteros tal que $\lvert A \rvert = 3 \, \text{m}$ para algún entero positivo m.

\textbf{Pregunta:} ¿Es posible particionar $A$ en m tríos disjuntos de forma que todos los tríos tengan la misma suma?

\medskip

Demostraremos que PT es NP-completo.

\subsection*{1. PT pertenece a NP}

Dada una partición candidata de A en m tríos, basta:

\begin{enumerate}
\item comprobar que cada elemento aparece exactamente una vez;
\item calcular la suma de cada trío;
\item verificar que todas las sumas coinciden.
\end{enumerate}

Todas estas operaciones pueden realizarse en tiempo polinómico respecto al tamaño de la entrada.

Por tanto, $PT \in NP$.


\subsection*{2. Problema NP-completo de partida}

Ahora que ya sabemos que el problema está en NP para demostrar que es completo nuestra intención será crear una reducción desde uno de los problemas ya conocidos y demostrados como NP-completos en clase. Por tanto, partiremnos del problema ACTRI (Acoplamiento por Tripletas). Recordemos en qué consistía dicho problema:

\medskip

\textbf{Instancia:}

\begin{itemize}
\item tres conjuntos disjuntos (W,X,Y) tales que

$|W|=|X|=|Y|=q;$

\item un conjunto de compatibilidades


$M\subseteq W\times X\times Y.$

\end{itemize}

\textbf{Pregunta:}

¿Existe un subconjunto $M'\subseteq M$ con exactamente q tripletas tal que cada elemento de $W\cup X\cup Y$ aparezca en una y solo una tripleta de $M'$?

\medskip

\subsection*{3. Observaciones y análisis previo a la reducción}

La elección de ACTRI como problema desde el que partir para realizar la reducción a priori parece bastante lógica pues tanto en ACTRI como en PT el objetivo es, de alguna forma, particionar un conjunto de objetos en tríos dijsuntos cumpliendo unas restricciones. No obstante, la diferencia clave entre ambos problemas se halla en cómo se definen esas restricciones sobre los tríos de elementos pues en ACTRI se hace a través del conjunto M de compatibilidades y en PT se consiguen mediante una asiganción de enteros a cada elemento (más bien cada elemento es un número entero) e imponiendo que la suma de los elementos de todos los tríos sea la misma. Por tanto, la dificultad en realizar la reducción yace en traducir de alguna forma las compatibilades impuestas en M a números enteros que se identifiquen con los elementos de los 3 conjuntos base de ACTRI y que a la vez cumplan la restricción antes mencionada de las sumas. Esta transformación no es para nada simple ni evidente y es por ello que, además de revisar al completo toda la teoría impartida en clase, he consultado el tan mencionado en la asignatura libro de Garey-Johnson. En él he llegado a ver que, entre muchos problemas tratados, se plantea uno llamado 3-PARTITION el cual, al leer su definición, nos damos cuenta de que se parece mucho al propio problema PT que queremos nosotros demostrar y en el libro se demuestra que 3-PARTITION es NP-completo partiendo del problema 3-DIMENSIONAL MATCHING que es el nombre que se le da exactamente al propio ACTRI. Por ello vamos a tomar la demostración de Garey-Johnson y adaptarla  para demostrar la NP-completitud de PT.


\subsection*{4. Esquema de las reducciones}

En el libro para llegar a demostrar que el problema 3-PARTITION es NP-completo primero desmuestra que otro problema auxiliar, 4-PARTITION, es NP-completo reduciéndolo desde el propio ACTRI. Cabe mencionar que en el libro, a diferencia de cómo hemos definido las reducciones imponiendo que estén en L, solo comprueba que sean polinómicas. En resumen, en el libro de Garey-Johnson se realizan las reducciones:

\vspace{1\baselineskip}
$ACTRI
\le_p
4\text{-PARTITION}
\le_p
3\text{-PARTITION}$

\vspace{1\baselineskip}
Definamos ahora los propios problemas 3-PARTITION Y 4-PARTITION que necesitaremos conocer para las reducciones:

\vspace{1\baselineskip}
\textbf{3-PARTITION}

\textbf{INSTANCIA:} Un conjunto finito $A$ de $3m$ elementos, una cota $B \in \mathbb{Z}^+$, y un ``tamaño'' $s(a) \in \mathbb{Z}^+$ para cada $a \in A$, tal que cada $s(a)$ satisface
\[
\frac{B}{4} < s(a) < \frac{B}{2},
\]
y además
\[
\sum_{a \in A} s(a) = mB.
\]

\textbf{PREGUNTA:} ¿Puede $A$ particionarse en $m$ subconjuntos disjuntos $S_1, S_2, \ldots, S_m$ tal que, para $1 \le i \le m$,
\[
\sum_{a \in S_i} s(a) = B?
\]
(Nótese que las restricciones sobre los tamaños de los elementos implican que cada $S_i$ debe contener exactamente tres elementos de $A$.)

\vspace{1\baselineskip}
\textbf{4-PARTITION}

\textbf{INSTANCIA:} Un conjunto finito $A$ de $4m$ elementos, una cota $B \in \mathbb{Z}^+$, y un ``tamaño'' $s(a) \in \mathbb{Z}^+$ para cada $a \in A$, tal que cada $s(a)$ satisface
\[
\frac{B}{5} < s(a) < \frac{B}{3},
\]
y además
\[
\sum_{a \in A} s(a) = mB.
\]

\textbf{PREGUNTA:} ¿Puede $A$ particionarse en $m$ subconjuntos disjuntos $S_1, S_2, \ldots, S_m$ tal que, para $1 \le i \le m$,
\[
\sum_{a \in S_i} s(a) = B?
\]
(Nótese que las restricciones sobre los tamaños de los elementos implican que cada $S_i$ debe contener exactamente cuatro elementos de $A$.)

\vspace{1\baselineskip}
Ya con estos dos problemas auxiliares definidos es fácil ver que, de forma análoga a cómo lo hicimos para PT, son ambos NP. Ya cerramos la estrategia que vamos a seguir: realizar las dos reducciones desde ACTRI hasta 3-PARTITION del Garey-Johnson y luego demostrar que el problema 3-PARTITION arriba definido es equivalente a nuestro problema de Prtición en Tríos, PT. Es decir:

$ACTRI
\propto
4\text{-PARTITION}
\propto
3\text{-PARTITION}
\equiv
PT.$

\subsection*{5. Reducción de ACTRI a 4-PARTITION}

Sea (W,X,Y,M) una instancia de ACTRI. Definimos r=32q. Para cada elemento $z\in W\cup X\cup Y$ que aparezca en N(z) tripletas de M, creamos los elementos $z[1],z[2],\ldots,z[N(z)]$. El elemento z[1] será la copia real y los restantes serán copias ficticias. A los elementos recién creados les asignamos los siguientes tamaños (donde $1 \le i, j, k \le q$ siempre):

\vspace{1\baselineskip}
$s(w_i[1])=10r^4+ir+1,$

$s(w_i[\ell])=11r^4+ir+1,
\qquad
2\le \ell \le N(w_i),$

$s(x_j[1])=10r^4+jr^2+2,$

$s(x_j[\ell])=11r^4+jr^2+2,
\qquad
2\le \ell \le N(x_j),$

$s(y_k[1])=10r^4+kr^3+4,$

$s(y_k[\ell])=8r^4+kr^3+4,
\qquad
2\le \ell \le N(y_k).$

\vspace{1\baselineskip}
Además, para cada tripleta $m=(w_i,x_j,y_k)\in M$ se crea un elemento $u_m$ con tamaño $s(u_m)
=
10r^4
-kr^3
-jr^2
-ir
+8.
$

\vspace{1\baselineskip}
Finalmente definimos $B=40r^4+15.$

\vspace{1\baselineskip}
Puede comprobarse inmediatamente que $s(u_m)+s(w_i[1])+s(x_j[1])+s(y_k[1])=B.$ Asimismo, si se utilizan simultáneamente las tres copias ficticias correspondientes, la suma vuelve a ser exactamente B. También no es complicado observar que el tamaño de cada elemento está estrictamente entre $B/3$ y $B/5$, y que la suma de todos los tamaños de los elementos es $|M|\,B$, como se requería. Además, observamos que el tamaño de cada elemento está acotado superiormente por $12r^4 \le 12\cdot 8^4\cdot |A|^4 < 2^{16}|A|^4$.

\subsection*{Equivalencia a través de la reducción}

\paragraph{Si ACTRI tiene solución.}

Sea $M'\subseteq M$ una solución de ACTRI. Para cada tripleta $m=(w_i,x_j,y_k)\in M'$ formamos el bloque $\{u_m,w_i[1],x_j[1],y_k[1]\}$. Para las tripletas restantes (es decir, para los $u_$ correspondientes a un $m\in M-M'$) utilizamos las copias ficticias correspondientes para formar grupos de 4 elementos.

Todos los bloques obtenidos tienen suma exactamente B, por lo que se obtiene una partición válida y, por tanto, una solución de 4-PARTITION.

\paragraph{Si existe una partición válida como solución de 4-PARTITION.}

Ahora supongamos que se nos da una 4-partición de la forma requerida. Consideremos cualquier tupla de esta 4-PARTITION. Considerando sucesivamente la suma de los tamaños de los elementos módulo $r$, $r^2$, $r^3$, $r^4$ y $r^5$, mostraremos que este conjunto de 4 elementos debe contener un elemento correspondiente a cada uno de los tresconjuntos originales de ACTRI, siendo las tres copias reales o las tres copias ficticias.

Primero, como $r > 4 \cdot 8 = 32$, sabemos que la suma módulo $r$ de los tamaños de los cuatro elementos, que debe ser igual a $B \pmod r = 15$, es la misma que la suma de los tamaños de los elementos cuando cada uno se toma módulo $r$ de antemano. La única manera de que esto sume $15$ es que el conjunto de 4 elementos contenga un elemento que corresponda a un miembro de $W$, uno que corresponda a un miembro de $X$, uno que corresponda a un miembro de $Y$, y uno que corresponda a un triple de $M$.

Llamemos $w_i, x_j, y_k$ a los miembros correspondientes de $W$, $X$ e $Y$, y llamemos $m_{i} = (w_{i'}, x_{j'}, y_{k'})$ al triple correspondiente de $M$ (el respectivo al $u_i$ del conjunto de 4 elementos que estamos considerando). Entonces la suma de los tamaños de los elementos módulo $r^2$ debe ser igual a $((i-i')r + 15) (\bmod r^2)$, y como $(i-i')r + 15 < r^2$, debemos tener

\[
B \pmod {r^2} = 15 - ((i-i')r + 15)
\]

Se sigue de esto que $i=i'$. De manera similar, como $(j-j')r^2 + 15 < r^3$, debemos tener

\[
B \pmod {r^3} = 15 - ((j-j')r^2 + 15)
\]

y por tanto $j=j'$, y como $(k-k')r^3 + 15 < r^4$, debemos tener

\[
B \pmod {r^4} = 15 - ((k-k')r^3 + 15)
\]

y por tanto $k=k'$. Así, $w_i, x_j$ y $y_k$ son efectivamente los tres miembros del triple $m_i$, y sabemos que el coeficiente de $r^4$ en la suma de los tamaños de los elementos es simplemente la suma de los coeficientes individuales de $r^4$. Nuestra elección de estos coeficientes en la construcción garantiza entonces que la única manera de que sumen $40$ es que los tres elementos sean todos ``reales'' o que los tres sean ``ficticios''.

La colección total de $3q$ elementos ``reales'', uno para cada miembro de $W \cup X \cup Y$, debe estar contenida en $q$ de nuestros 4-conjuntos dados, cada uno de estos 4-conjuntos consistente en un elemento correspondiente a un triple de $M$ y los tres elementos ``reales'' correspondientes a los miembros de ese triple. Esos $q$ triples de $M$ proporcionan la correspondencia deseada.

\subsection*{Complejidad de la reducción}
Por último en este paso nos queda comprobar que la reducción recién hecha está en $L$, es decir, que tiene compplejidad espacio logarítmica.

Sea I=(W,X,Y,M) una instancia de ACTRI.

La construcción de la instancia de 4-PARTITION consiste únicamente en generar una colección de elementos cuyos tamaños vienen dados por expresiones explícitas en función de:
$(
q,\quad i,\quad j,\quad k,\quad N(z).
)$

Para producir cada elemento de la salida no es necesario almacenar simultáneamente toda la instancia construida. Basta recorrer secuencialmente la entrada y mantener:

\begin{itemize}
\item los índices actuales $(i,j,k)$;
\item el valor de $q$;
\item contadores para las ocurrencias $N(z)$;
\item algunas variables auxiliares necesarias para evaluar las expresiones que definen los tamaños.
\end{itemize}

Todos estos valores están acotados por un polinomio en el tamaño de la entrada, por lo que pueden almacenarse utilizando $O(\log |I|)$ bits.

Además, los tamaños generados son expresiones del tipo $10r^4+ir+1,
\qquad
10r^4-jr^2-ir+8,$ donde $r=32q$.

Dichos valores pueden calcularse mediante operaciones aritméticas sobre enteros de longitud $O(\log |I|)$, reutilizando el mismo espacio de trabajo.

Por tanto, la instancia de salida puede generarse secuencialmente utilizando únicamente espacio logarítmico.

En consecuencia, la transformación descrita constituye una reducción en espacio logarítmico estando por tanto en $L$ y pudiendo afirmar finalmente que 

\vspace{1\baselineskip}
$ACTRI \propto 4\text{-PARTITION}$

\vspace{1\baselineskip}
Por lo tanto, ya hemos demostrado que 4-PARTITION es NP-completo.

\subsection*{6. Reducción de 4-PARTITION a 3-PARTITION}

Sea $A=\{a_1,a_2,\ldots,a_{4n}\}$, una cota $B$, y tamaños de elementos $s(a)$ que satisfacen $B/5 < s(a) < B/3$ y $s(a) \le 2^{16}\,|A|^4$ una especificación de cualquier instancia de 4-PARTITION de este tipo. Nuestra instancia correspondiente de 3-PARTITION tendrá $24n^2-3n$ elementos, uno por cada elemento de $A$, dos por cada par de elementos de $A$, y $8n^2-3n$ elementos ``relleno''.

A cada elemento $a_i \in A$ le corresponde un elemento ``regular'' $w_i$, con tamaño definido por

\[
s'(w_i)=4\cdot(5B+s(a_i))+1.
\]

donde usamos $s'(\cdot)$ para denotar la función de tamaño en nuestra instancia de 3-PARTITION. A cada par de elementos $a_i,a_j\in A$ le corresponden dos elementos ``emparejamiento'', $u[i,j]$ y $\bar{u}[i,j]$, con tamaños definidos por

\[
s'(u[i,j])=4\cdot(6B-s(a_i)-s(a_j))+2,
\]

\[
s'(\bar{u}[i,j])=4\cdot(5B+s(a_i)+s(a_j))+2.
\]

Finalmente, para $1 \le k \le 8n^2-3n$, tenemos un elemento ``relleno'' $u_k^*$ con tamaño

\[
s'(u_k^*)=20B.
\]

La cota $B'$ para nuestra instancia de 3-PARTITION será $64B+4$.

Es fácil ver que el tamaño de cada elemento está estrictamente entre $B'/4=16B+1$ y $B'/2=32B+2$, y que la suma de todos los tamaños de los elementos es igual a $(8n^2-n)B'$. Además, como los elementos de $A$ están restringidos a tener tamaños no mayores que $2^{16}|A|^4$, los tamaños en la instancia de 3-PARTITION también satisfarán una cota polinómica en términos de $|A|$, de ahí en términos del número de elementos en la instancia construida $I'$, en la longitud $|I'|$. Debemos ver ahora la equivalencia a través de esta reducción entre ambos problemas, es decir, que existe una 3-PARTITION para la instancia construida si y solo si existe una 4-PARTITITON para la instancia original.


\subsection*{Equivalencia a través de la reducción}

\paragraph{Si existe una partición válida como solución de 4-PARTITION.}

Primero, supongamos que tenemos una instancia de 4-PARTITION. La correspondiente 3-PARTITION se construye así: dividimos arbitrariamente cada conjunto de la 4-PARTITION $\{a_i,a_j,a_k,a_l\}$ en dos conjuntos de dos elementos, por ejemplo $\{a_i,a_j\}$ y $\{a_k,a_l\}$. Nuestra 3-PARTITION contendrá entonces los dos conjuntos de 3 elementos $\{w_i,w_j,u[i,j]\}$ y $\{w_k,w_l,\bar{u}[k,l]\}$.

Los tamaños de los elementos en cada uno de estas dos tripletas creadas suman $B'$, ya que $s(a_i)+s(a_j)+s(a_k)+s(a_l)=B$. Haciendo esto para cada uno de los $q$ conjuntos de la 4-PARTITION de partida, obtenemos $2n$ 3-conjuntos que contienen a todos los elementos ``reales'' y $n$ pares emparejados de elementos ``emparejamiento''. Esto deja $8n^2-3n$ pares emparejados de elementos ``emparejamiento'' y $8n^2-3n$ elementos ``relleno''. Como la suma de los tamaños de dos elementos ``emparejamiento'' emparejados es $44B+4=B'-20B$, cada par emparejado puede agruparse con uno de los elementos ``relleno'' para completar la 3-partición deseada.

\paragraph{Si existe una partición válida como solución de 3-PARTITION.}

Ahora supongamos que se nos da una instancia de 3-PARTITION creada a partir de la construcción de la reducción. Considerando los tamaños de los elementos módulo $4$, vemos que ninguna tripleta puede contener un número impar de elementos ``reales'' o ``regulares'', por lo que una tripleta solo puede contener tres elementos ``emparejamiento'' y un elemento ``relleno'' o dos elementos ``reales'' y un elemento ``relleno''. Se deduce que la 3-PARTITION dada está formada por $2n$ tripletas que contienen dos elementos ``reales'' y un elemento ``emparejamiento'', y por $8n^2-3n$ tripletas que contienen dos elementos ``emparejamiento'' y un elemento ``relleno''.

Consideremos cualquiera de los segundos $8n^2-3n$ tipos de tripletas; llamemos $u[i,j]$ (o $\bar{u}[k,l]$) a uno de los dos elementos ``emparejamiento'' de ese conjunto. Si el otro elemento ``emparejamiento'' de ese 3-conjunto no es $\bar{u}[i,j]$ (o $u[k,l]$), entonces debe tener el mismo tamaño que ese elemento emparejado y pueden intercambiarse para obtener una 3-PARTITION equivalente. Esta operación puede repetirse hasta que obtengamos una 3-PARTITION equivalente e igual de válida en la que cada elemento ``relleno'' aparezca junto con un par emparejado $u[i,j],\bar{u}[i,j]$. Así, cualquier elemento ``emparejamiento'' que aparezca con dos elementos ``reales'' en esta 3-PARTITION debe hacerlo de forma que también ocurra una ``correspondencia'' en esa tripleta. Esto divide los tríos que contienen elementos ``reales'' en pares de tripletas.

Como los dos elementos ``emparejamiento'' de cada par de tripletas están emparejados, sus tamaños suman $44B+4$, y por tanto los tamaños de los cuatro elementos ``reales'' deben sumar $84B+4$. Esto implica que los cuatro elementos correspondientes de $A$ forman un 4-conjunto cuyo tamaño es $B$. Por tanto, esos $n$ pares de 3-conjuntos proporcionan la 4-PARTITION requerida.

\subsection*{Complejidad de la reducción}
Por último en este paso nos queda comprobar que la nueva reducción recién hecha también está en $L$.

\vspace{1\baselineskip}
Sea $A=\{a_1,a_2,\ldots,a_{4n}\}$ una instancia de 4-PARTITION.

La construcción de la instancia correspondiente de 3-PARTITION consiste en generar:

\begin{itemize}
\item un elemento regular $w_i$ para cada elemento $a_i$;
\item dos elementos de emparejamiento $u[i,j]$ y $\bar{u}[i,j]$ para cada par de elementos $a_i,a_j$;
\item $8n^2-3n$ elementos de relleno.
\end{itemize}

Para producir cualquiera de estos elementos no es necesario almacenar simultáneamente toda la instancia construida. Basta recorrer secuencialmente la entrada manteniendo:

\begin{itemize}
\item los índices actuales $i$ y $j$;
\item el valor de $n$;
\item los tamaños $s(a_i)$ y $s(a_j)$ de los elementos que intervienen en el cálculo;
\item algunas variables auxiliares necesarias para efectuar las operaciones aritméticas.
\end{itemize}

Todos estos valores están acotados por un polinomio en el tamaño de la entrada, por lo que pueden representarse utilizando únicamente
$O(\log |I|)$
bits.

Además, los tamaños generados vienen dados por expresiones explícitas:

\[
s'(w_i)=4(5B+s(a_i))+1,
\]

\[
s'(u[i,j])=4(6B-s(a_i)-s(a_j))+2,
\]

\[
s'(\bar{u}[i,j])=4(5B+s(a_i)+s(a_j))+2,
\]

\[
s'(u_k^*)=20B,
\]

las cuales pueden calcularse mediante operaciones aritméticas sobre enteros de longitud polinómica reutilizando siempre la misma memoria de trabajo.

Por tanto, la instancia de salida puede generarse secuencialmente utilizando únicamente espacio logarítmico respecto al tamaño de la entrada.

\vspace{1\baselineskip}
En consecuencia, la transformación descrita constituye una reducción en espacio logarítmico estando por tanto en $L$ y pudiendo afirmar finalmente que 

\vspace{1\baselineskip}
$4\text{-PARTITION} \propto 3\text{-PARTITION}$

\vspace{1\baselineskip}
Por lo tanto, ya hemos demostrado que 3-PARTITITON es NP-completo.

\subsection*{7. Equivalencia entre 3-PARTITION y PT}

Por último nos queda ver que nuestro problema original, PT, es equivalente al problema 3-PARTITION o que al menos lo es al partir de ACTRI a través de las 2 reducciones construidas.


\vspace{1\baselineskip}
Sea $(A_3,B)$ una instancia de 3-PARTITION. La transformación para obtener una instancia de PT consiste simplemente en conservar el conjunto de elementos $A_3$ pero considerando como el $A$ de PT el conjunto $A = \{s(a) : a \in A_3\}$. Nos surge la pregunta clave: ¿$|A|=|A_3|$? Pues para que ello pase la función $s$ debería ser inyectiva pero en la definición propia del problema 3-PARTITION aparentemente no lo pide estrictamente. No obstante, si consideramos que realmente la isntancia de 3-PARTITION que estamos considerando proviene de concatenar las dos reducciones previas desde una instancia de ACTRI es fácil ver que por las construcciones que dimos en ambas reducciones en este caso la función $s$ sí que es efectivamente inyectiva. Por lo tanto, sí que podemos establecer una biyección completa entre $A$ y $A_3$. 


\vspace{1\baselineskip}
Así, supongamos que $|A_3|=3m$ y que $\sum_{a\in A_3} s(a) = mB$.

Si existe una solución de 3-PARTITION, las m tripletas obtenidas suman todas exactamente B, luego los tríos obtenidos a través de la biyección entre $A_3$ y $A$ constituyen una solución de PT.


\vspace{1\baselineskip}
Recíprocamente, si PT admite una partición en tríos de suma común K, entonces $mK=\sum_{a\in A}$. Tomando $A_3 := A$, $B := K$  y como $s : A_3 \to \mathbb{Z}^+$ la identidad en $\mathbb{Z}$ restringida a $A_3$ tengo una solución de 3-PARTITION puesto que la otra condición de este problema impuesta sobre la aplicación $s$, la de que $\frac{B}{4} < s(a) < \frac{B}{2}$, es trivial que se cumple para todos los elementos  de una solución de PT pues de no hacerlo, como todos los tríos deben sumar B pues es solución, entonces

\vspace{1\baselineskip}
$\forall a \in A : a \le \frac{B}{4} \implies \exists b \in A : b \quad está \quad en \quad el \quad mismo \quad trío \quad que \quad a \quad y \quad b \ge \frac{B}{2}.$ 

\vspace{1\baselineskip}
Es claro que en esa situación podemos restarle a $b$ lo que "le sobra" para no quedarse por encima de $\frac{B}{2}$ y sumárselo a $a$ para compensar. Si aún así sigue ocurriendo que $a \le \frac{B}{4}$ será porque el tercer elemento del trío, $c$, le sucede también $c \ge \frac{B}{2}$ y, por tanto, podemos repetir lo hecho antes con b. En esta situación seguro que ya los 3 valores del trío deben estar dentro del intervalo $]\frac{B}{4}, \frac{B}{2}[$.

De igual forma se sigue un razonamiento análogo $\forall a \in A : a \ge \frac{B}{2}$.

Por tanto, es claro que podíamos suponer sin pérdida de generalidad que 

$ \quad \quad \quad \quad \quad \quad \quad \quad \quad \quad\frac{B}{4} < s(a) < \frac{B}{2} \quad \forall a \in A$.

\vspace{1\baselineskip}
Luego los tríos de PT constituyen también una solución de 3-PARTITION.

Así, tenemos al fin que $3\text{-PARTITION}
\equiv
PT$.

\vspace{1\baselineskip}
Podemos pasar ya a una recapitulación de todo lo conseguido y obtener la conclusión final.

\subsection*{7. Conclusión}

Valiéndonos de los conocimientos de la asignatura y del libro de referencia de Garey-Johnson hemos demostrado que:

\vspace{1\baselineskip}
$ \quad \quad \quad \quad \quad \quad \quad ACTRI
\propto 
4\text{-PARTITION}
\propto
3\text{-PARTITION}
\equiv
PT$

\vspace{1\baselineskip}
Así, como ACTRI es NP-completo y PT pertenece a NP, se concluye que PT es NP-completo:

\vspace{1\baselineskip}
$ \quad \quad \quad \quad \quad \quad \quad \quad \boxed{\text{Partición en tríos es NP-completo}.}$


\end{document}
